\chapter{Abdominal Aortic Aneurysm Repair}

\section*{What Is This Surgery?}
Abdominal aortic aneurysm (AAA) repair is a major open surgical procedure designed to replace a pathologically dilated and weakened segment of the abdominal aorta with a synthetic graft. This intervention is primarily performed to prevent the catastrophic and often fatal event of aneurysm rupture, which involves massive internal hemorrhage. Typically, the procedure addresses aneurysms located in the infrarenal aorta, though juxtarenal or suprarenal extensions may necessitate more complex techniques. While endovascular aneurysm repair (EVAR) has become a common alternative, open repair remains a critical and definitive treatment, particularly for patients with complex anatomy, failed EVAR, or those presenting with rupture. It offers a durable, long-term solution by physically replacing the diseased aortic segment.

\section*{Alternative Names}
\begin{itemize}
  \item Open Abdominal Aortic Aneurysmectomy
  \item Aortic Graft Placement for Aneurysm
  \item Aneurysmectomy with Graft Interposition
  \item Conventional AAA Repair
  \item Infrarenal Aortic Aneurysm Repair
  \item Aortic Bypass for Aneurysm
  \item Aortoiliac Bypass (if bifurcated graft used)
\end{itemize}

\section*{Relevant Surgical Anatomy}
The abdominal aorta is the largest artery in the abdomen, originating at the diaphragm and typically bifurcating into the common iliac arteries at the level of the fourth lumbar vertebra (L4). An abdominal aortic aneurysm most commonly occurs in the infrarenal segment, meaning below the origin of the renal arteries. Key anatomical structures relevant to open AAA repair include:

\begin{itemize}
  \item **Abdominal Aorta:** The primary vessel of concern, typically atherosclerotic and dilated. Its anterior surface is covered by the posterior peritoneum, and it lies anterior to the lumbar vertebrae.
  \item **Renal Arteries:** Arise laterally from the aorta, typically at the level of L1-L2. The proximal anastomosis of the graft is usually performed just inferior to these vessels to preserve renal perfusion.
  \item **Inferior Mesenteric Artery (IMA):** Originates from the anterior aspect of the aorta, typically 3-5 cm above the aortic bifurcation. It supplies the distal colon. Its preservation or revascularization is crucial if collateral flow is compromised.
  \item **Lumbar Arteries:** Paired arteries arising from the posterior aspect of the aorta, supplying the posterior abdominal wall and spinal cord. These are typically ligated or oversewn from within the aneurysm sac.
  \item **Common Iliac Arteries:** The terminal branches of the aorta, which further divide into internal and external iliac arteries. Distal anastomoses for bifurcated grafts are typically made to the common iliac arteries.
  \item **Left Renal Vein:** Crosses anterior to the aorta, typically just below the superior mesenteric artery and above the renal arteries. It must be carefully mobilized and protected during proximal aortic dissection and clamping.
  \item **Duodenum:** The third part of the duodenum crosses anterior to the aorta at the level of L3. It is mobilized superiorly and to the right to expose the infrarenal aorta.
  \item **Ureters:** Lie in the retroperitoneum, lateral to the great vessels. They must be identified and protected from injury during dissection, especially in cases of inflammatory aneurysms.
  \item **Sympathetic Nerves:** Plexuses of sympathetic nerves lie anterior and lateral to the aorta, particularly around the bifurcation. Injury to these nerves can lead to ejaculatory dysfunction or impotence.
\end{itemize}
The retroperitoneal space, where the aorta resides, is accessed via a transperitoneal approach, requiring careful mobilization of overlying viscera.

\section*{Surgical Principle \& Rationale}
The fundamental principle of open abdominal aortic aneurysm repair is the definitive exclusion of the diseased, aneurysmal segment of the aorta from the systemic circulation and the restoration of normal aortic continuity and blood flow using a prosthetic conduit. This is achieved by surgically exposing the aorta, temporarily clamping it above and below the aneurysm to control blood flow, and then incising the aneurysm sac longitudinally. Any intraluminal thrombus is meticulously removed. A synthetic graft, typically made of woven or knitted Dacron (polyethylene terephthalate) or PTFE (polytetrafluoroethylene), is then precisely sewn into place, creating a new, strong conduit for blood flow. The proximal anastomosis is usually performed just below the renal arteries, and the distal anastomoses are made to the common iliac arteries (for a bifurcated graft) or to the infrarenal aorta (for a tube graft). The native aneurysm sac is then typically closed around the graft to protect it from surrounding structures and prevent the formation of an aortoenteric fistula. This technique effectively eliminates the risk of rupture from the original weakened aortic wall.

The rationale for this approach is to provide a durable, low-resistance pathway for blood flow while completely removing the risk posed by the weakened aortic wall. Technical goals include achieving secure, hemostatic anastomoses that can withstand systemic arterial pressure, maintaining adequate perfusion to vital organs during the period of aortic clamping, and ensuring the long-term patency of the graft. The choice between a tube graft (for aneurysms confined to the aorta) and a bifurcated graft (for aneurysms extending into the iliac arteries) depends on the specific anatomical extent of the disease, aiming to restore normal flow to both lower extremities.

\section*{History \& Surgical Evolution}
The concept of surgically treating aortic aneurysms dates back to the early 20th century, with initial attempts involving ligation or wrapping of the aneurysm, often with poor outcomes. A significant breakthrough occurred in 1951 when Dr. Charles Dubost in Paris performed the first successful resection of an abdominal aortic aneurysm and replaced it with a homograft (cadaveric aorta). This pioneering work demonstrated the feasibility of direct surgical repair. The subsequent development of synthetic graft materials, notably Dacron by Dr. Arthur Voorhees and Dr. Arthur Blakemore in the mid-1950s, revolutionized vascular surgery by providing readily available, durable, and biocompatible prostheses, eliminating the limitations and supply issues of homografts.

Throughout the latter half of the 20th century, open AAA repair evolved into a standardized and highly effective procedure, with refinements in surgical technique, anesthetic management, and perioperative care significantly improving patient outcomes. Key advancements included improved understanding of aortic anatomy, development of specialized vascular instruments, and better methods for intraoperative blood management. It remained the gold standard for AAA treatment until the advent of endovascular aneurysm repair (EVAR) in the 1990s, pioneered by Dr. Juan Parodi. While EVAR offered a less invasive alternative, open repair continues to be a vital procedure, particularly for complex cases, ruptured aneurysms, or when endovascular options are unsuitable due to anatomical constraints or prior EVAR failure.

\section{*{Clinical Indications}}
Open abdominal aortic aneurysm repair is indicated for specific clinical scenarios to prevent the life-threatening complication of aneurysm rupture or to address existing complications. Key reasons include:
\begin{itemize}
  \item Aneurysm size: Typically for infrarenal or juxtarenal AAA measuring 5.5 cm or larger in men, or 5.0 cm or larger in women, due to the increased risk of rupture at these diameters.
  \item Rapid aneurysm expansion: An increase in aneurysm diameter of 0.5 cm or more within a 6-month period, regardless of absolute size, indicating an unstable aneurysm.
  \item Symptomatic aneurysm: Presence of new or worsening abdominal, back, or flank pain, which may indicate impending rupture, rapid expansion, or contained rupture.
  \item Aneurysm rupture: An emergent, life-saving procedure performed immediately when an aneurysm has ruptured, presenting with severe pain, hypotension, and shock.
  \item Distal embolization: Evidence of thrombus within the aneurysm sac causing emboli to the lower extremities or other organs, leading to acute limb ischemia or organ dysfunction.
  \item Anatomical unsuitability for EVAR: When the aneurysm morphology (e.g., short or angulated neck, severe iliac tortuosity, extensive thrombus, or large diameter) precludes safe and effective endovascular repair.
  \item Infected (mycotic) aneurysm: Open repair with debridement and often extra-anatomic bypass is typically preferred for infected aneurysms due to the high risk of graft infection with endovascular approaches.
  \item Failed previous EVAR: For patients who develop complications after EVAR, such as persistent or enlarging endoleak, graft infection, device migration, or rupture, requiring conversion to open repair.
\end{itemize}

\section*{Clinical Positioning}
Open abdominal aortic aneurysm repair serves as both a primary and a secondary definitive treatment for AAA, depending on the clinical context and patient factors. Historically, it was the primary and only definitive surgical treatment for decades. With the advent of endovascular aneurysm repair (EVAR) in the 1990s, open repair's role has evolved, but it remains indispensable.

As a primary procedure, open repair is chosen when a patient meets the size or symptomatic criteria for intervention and is a suitable surgical candidate, but their aneurysm anatomy is not amenable to EVAR. This includes cases with a very short or absent infrarenal neck, severe neck angulation, extensive iliac artery disease, significant thrombus burden, or severe calcification that would compromise EVAR success. It is also the primary choice for infected aneurysms, those with connective tissue disorders (e.g., Marfan syndrome) that may predispose to graft complications with EVAR, or in younger, healthier patients where the long-term durability of open repair is favored over potential re-interventions associated with EVAR.

As a secondary procedure, open repair is indicated for patients who have previously undergone EVAR but have developed complications such as persistent or enlarging endoleaks (Type I or III), graft infection, device migration, or rupture after EVAR. In these "failed EVAR" scenarios, open conversion is often necessary to address the ongoing pathology, which can be technically more challenging due to previous scarring and the presence of the endograft. Furthermore, in emergent situations like ruptured AAA, open repair remains a critical option, especially if the patient's condition or anatomy makes EVAR technically challenging or time-prohibitive, or if EVAR equipment is not immediately available.

\section*{Surgical Technique \& Procedure}
Open abdominal aortic aneurysm repair is a complex surgical procedure performed under general anesthesia, often supplemented with epidural anesthesia for postoperative pain control.

The patient is positioned supine on the operating table. A long midline incision is made from the xiphoid process to the pubis, or sometimes a transverse incision, to gain access to the abdominal cavity. The small bowel is typically retracted to the right side of the abdomen using a self-retaining retractor system, exposing the retroperitoneum. The posterior peritoneum overlying the aorta is then incised longitudinally, and the third part of the duodenum is mobilized superiorly and to the right to expose the infrarenal aorta. The aorta is carefully dissected free from surrounding tissues, including the left renal vein, inferior mesenteric artery, and lumbar arteries, ensuring adequate length for clamping. Control of the aorta is achieved by placing vascular clamps: one proximally, typically just below the renal arteries, and two distally, usually on the common iliac arteries. Heparin is administered intravenously before clamping to prevent thrombus formation.

Once the aorta is clamped, the aneurysm sac is incised longitudinally. Any intraluminal thrombus and atherosclerotic debris are meticulously removed. The lumbar arteries originating from the aneurysm sac are identified and ligated or oversewn from within to prevent back-bleeding. The inferior mesenteric artery (IMA) ostium is assessed; if it is patent and there is concern for colonic ischemia, it may be reimplanted into the graft. A synthetic graft, either a tube graft (for aneurysms confined to the aorta) or a bifurcated graft (for aneurysms extending into the iliac arteries), is selected based on the aneurysm's extent. The proximal anastomosis is performed first, sewing the graft to the healthy aorta just below the renal arteries using a continuous suture technique (e.g., 3-0 or 4-0 polypropylene). The distal anastomoses are then completed, connecting the graft limbs to the common iliac arteries (for a bifurcated graft) or the distal aorta (for a tube graft). Before completing the final anastomosis, the clamps are briefly released to flush out any air or debris. Once all anastomoses are complete and secure, the clamps are removed, restoring blood flow through the new graft. Hemostasis is meticulously achieved. The native aneurysm sac is then closed over the graft (aneurysmorrhaphy) to provide a protective layer and prevent aortoenteric fistula formation. The posterior peritoneum is closed, and the abdominal wall is closed in layers, often with retention sutures in high-risk patients. Drains may be placed in the retroperitoneum or abdominal cavity before closure.

\section*{Preoperative Workup \& Preparation}
Extensive preoperative preparation is crucial for patients undergoing open abdominal aortic aneurysm repair to optimize their condition and minimize risks. This comprehensive evaluation aims to identify and mitigate cardiac, pulmonary, and renal comorbidities.

\begin{itemize}
  \item **Comprehensive Medical Evaluation:**
    \begin{itemize}
      \item **Cardiac Assessment:** Electrocardiogram (ECG), echocardiogram, and often a cardiac stress test (pharmacologic or exercise) to assess myocardial function and identify coronary artery disease. Patients with significant cardiac risk may require coronary revascularization prior to AAA repair, following guidelines from the American College of Cardiology/American Heart Association.
      \item **Pulmonary Assessment:** Pulmonary function tests (PFTs) for patients with a history of smoking or lung disease. Smoking cessation is strongly encouraged for several weeks preoperatively to improve lung function and reduce postoperative pulmonary complications.
      \item **Renal Assessment:** Serum creatinine, blood urea nitrogen (BUN), and estimated glomerular filtration rate (eGFR) to assess kidney function, as renal complications are a significant risk, especially with aortic clamping.
      \item **Hematologic Assessment:** Complete blood count (CBC), coagulation profile (PT, PTT, INR), and blood type and cross-match for several units of blood, as significant blood loss is possible.
    \end{itemize}
  \item **Imaging:** Computed tomography angiography (CTA) of the abdomen and pelvis with intravenous contrast and 3D reconstruction is essential to precisely map the aneurysm's anatomy, its relation to renal and visceral arteries, the extent of iliac involvement, and the presence of thrombus or calcification.
  \item **Medication Management:**
    \begin{itemize}
      \item Beta-blockers are often initiated or continued to reduce cardiac risk and control heart rate.
      \item Anticoagulants and antiplatelet medications (e.g., aspirin, clopidogrel) are typically held for 5-7 days prior to surgery, as directed by the surgeon and cardiologist, to minimize bleeding risk.
      \item Diabetes management: Blood glucose levels must be tightly controlled preoperatively and perioperatively.
    \end{itemize}
  \item **Bowel Preparation:** While not universally mandated, some surgeons may order a mechanical bowel preparation (e.g., polyethylene glycol solution) to reduce the risk of infection if the bowel is inadvertently injured during dissection.
  \item **NPO Status:** Patients must be NPO (nil per os) for at least 6-8 hours prior to surgery, as per anesthesia protocol, to prevent aspiration.
  \item **Antibiotic Prophylaxis:** Intravenous broad-spectrum antibiotics (e.g., cefazolin) are administered shortly before incision to reduce the risk of surgical site infection and devastating graft infection.
  \item **Informed Consent:** Detailed discussion with the patient and family regarding the procedure, its risks, benefits, alternatives (including EVAR and surveillance), and potential long-term implications, obtaining fully informed consent.
\end{itemize}

\section*{Contraindications \& Precautions}
Open abdominal aortic aneurysm repair is a major surgery with significant physiological demands, necessitating careful patient selection.

**Absolute Contraindications:**
\begin{itemize}
  \item **Active Systemic Infection or Sepsis:** The presence of an active infection, especially bacteremia, significantly increases the risk of graft infection, a devastating complication that often requires graft removal.
  \item **Severe, Uncorrectable Coagulopathy:** Uncontrolled bleeding disorders (e.g., severe hemophilia, severe thrombocytopenia) make the procedure prohibitively risky due to the high potential for hemorrhage.
  \item **Prohibitive Surgical Risk:** Patients with end-stage organ failure (e.g., severe heart failure with ejection fraction <20\%, recent myocardial infarction within 3-6 months, severe uncompensated respiratory failure requiring chronic ventilation, end-stage liver disease with coagulopathy) where the risk of mortality or severe morbidity from surgery far outweighs the benefit of aneurysm repair.
  \item **Limited Life Expectancy:** Patients with a very short life expectancy (e.g., less than 1-2 years) due to other terminal illnesses (e.g., metastatic cancer, advanced dementia), where the benefits of elective AAA repair are unlikely to be realized.
\end{itemize}

**Relative Contraindications and Precautions:**
\begin{itemize}
  \item **Severe Cardiopulmonary Disease:** While not always absolute, severe chronic obstructive pulmonary disease (COPD) with FEV1 <1 liter, unstable angina, or significant left ventricular dysfunction may increase perioperative risk. These patients require intensive optimization and may be better candidates for EVAR or surveillance if appropriate.
  \item **Extreme Age or Frailty:** Very elderly or frail patients may have diminished physiological reserve, making recovery from major surgery challenging. A comprehensive geriatric assessment is often warranted, focusing on functional status and cognitive reserve.
  \item **Hostile Abdomen:** Multiple previous abdominal surgeries leading to dense adhesions can make surgical access difficult, increase operative time, and raise the risk of bowel injury.
  \item **Morbid Obesity:** Can complicate surgical exposure, increase operative time, and elevate risks of wound complications, respiratory issues, and deep vein thrombosis.
  \item **Renal Insufficiency:** Patients with pre-existing chronic kidney disease are at higher risk for acute kidney injury postoperatively, especially with suprarenal clamping. Careful hydration, avoidance of nephrotoxic agents, and consideration of renal protection strategies are crucial.
  \item **Connective Tissue Disorders:** Conditions like Marfan syndrome or Ehlers-Danlos syndrome can lead to fragile aortic tissue, making anastomoses technically challenging and increasing the risk of pseudoaneurysm formation or aortic dissection.
  \item **Inflammatory Aneurysm:** While not a contraindication, inflammatory aneurysms require careful dissection due to dense perianeurysmal fibrosis, which can encase and obscure vital structures like the ureters or duodenum, increasing the risk of injury.
\item **Horseshoe Kidney:** This congenital anomaly can complicate proximal aortic clamping and dissection, requiring careful preoperative planning and intraoperative identification.
\end{itemize}

\section*{Complications \& Risks}
Open abdominal aortic aneurysm repair carries significant risks, both general surgical risks and those specific to aortic surgery, reflecting its complexity and the patient's underlying cardiovascular disease.

\begin{itemize}
  \item **General Surgical Risks:**
    \begin{itemize}
      \item **Bleeding:** Significant intraoperative or postoperative hemorrhage requiring blood transfusions or reoperation. This is a major concern given the size of the vessels involved.
      \item **Infection:** Surgical site infection (wound infection), pneumonia, urinary tract infection, or, rarely, severe graft infection (0.5-2\%), which is a devastating complication.
      \item **Anesthesia Complications:** Myocardial infarction (heart attack), stroke, adverse drug reactions, aspiration pneumonia, deep vein thrombosis (DVT) leading to pulmonary embolism (PE).
      \item **Pain:** Significant postoperative pain requiring intensive analgesia.
    \end{itemize}
  \item **Procedure-Specific Risks:**
    \begin{itemize}
      \item **Renal Failure:** Acute kidney injury (AKI) is a significant risk (5-20\%), especially if suprarenal clamping is required, in patients with pre-existing renal insufficiency, or due to prolonged hypotension.
      \item **Bowel Ischemia:** Insufficient blood supply to the colon (colonic ischemia), particularly the left colon, due to ligation of the inferior mesenteric artery or interruption of collateral flow. This can lead to necrosis, requiring colectomy, and carries a high mortality rate.
      \item **Spinal Cord Ischemia:** Although rare (0.2-2\%), interruption of blood flow to the spinal cord, especially with extensive aortic clamping, sacrifice of critical lumbar arteries, or prolonged hypotension, can result in paraplegia or paraparesis.
      \item **Graft Infection:** A rare but devastating complication (0.5-2\%) that can lead to sepsis, anastomotic pseudoaneurysm, or aortoenteric fistula, often requiring graft removal and extra-anatomic bypass.
      \item **Anastomotic Pseudoaneurysm:** Weakness at the suture line between the graft and native artery, leading to a contained rupture that can expand over time, requiring re-intervention.
      \item **Distal Embolization:** Release of plaque or thrombus from the aneurysm or clamped aorta, causing blockages in the arteries of the legs, kidneys, or other organs, leading to ischemia.
      \item **Lymphatic Leak/Lymphocele:** Injury to lymphatic channels during dissection, particularly in the groin or retroperitoneum, can lead to persistent drainage or fluid collection.
      \item **Erectile Dysfunction/Ejaculatory Dysfunction:** Injury to the sympathetic nerves during dissection, particularly around the aortic bifurcation, can affect ejaculatory function and potency.
      \item **Incisional Hernia:** Weakness in the abdominal wall closure can lead to a hernia at the incision site, requiring subsequent repair in up to 10-20\% of patients.
      \item **Aortoenteric Fistula:** A rare but life-threatening connection between the aortic graft and the gastrointestinal tract, typically the duodenum, leading to massive gastrointestinal bleeding. This is often a late complication of graft infection.
    \end{itemize}
\end{itemize}

\section*{Postoperative Recovery Timeline}
Immediately after open abdominal aortic aneurysm repair, the patient is transferred to the intensive care unit (ICU) for close monitoring. This includes continuous monitoring of vital signs, cardiac rhythm, urine output, central venous pressure, and arterial blood pressure. Mechanical ventilation is often required initially, with extubation typically occurring within 12-24 hours, depending on the patient's respiratory status and overall stability. Pain management is a priority, often utilizing epidural analgesia, patient-controlled analgesia (PCA) with opioids, or a combination of methods to ensure comfort and facilitate early mobilization. Fluid balance is carefully managed to maintain adequate hydration and renal perfusion, and blood transfusions may be necessary. Early ambulation is encouraged as soon as the patient is stable and able, usually within 24-48 hours, to prevent complications like deep vein thrombosis, pulmonary embolism, and atelectasis.

Over the next several days, the patient's condition is progressively monitored and managed. Bowel function, which is often temporarily impaired (ileus) due to abdominal surgery, is closely observed; a clear liquid diet is typically started once bowel sounds return and flatus is passed, advancing to a regular diet as tolerated. Surgical drains are usually removed when drainage is minimal, typically within 2-5 days. The hospital stay for open AAA repair is typically longer than for EVAR, ranging from 4 to 10 days, sometimes longer depending on recovery and complications. Upon discharge, patients are advised on wound care, activity restrictions (avoiding heavy lifting, straining, or strenuous exercise for 6-8 weeks to allow abdominal wall healing), and medication management. Full recovery can take several weeks to months, with gradual return to normal activities, and patients often experience fatigue for an extended period.

\section*{Surgical Findings \& Pathology}
During open abdominal aortic aneurysm repair, the surgeon documents several key findings. These typically include the precise location and extent of the aneurysm (e.g., infrarenal, juxtarenal, or suprarenal extension), its maximum diameter, the presence and amount of intraluminal thrombus, the degree of calcification in the aortic wall, and the involvement of the common iliac arteries. The quality of the aortic neck (proximal landing zone) and distal vessels (iliac arteries) for anastomosis is also noted. Any inflammatory changes around the aneurysm, adherence to adjacent structures (like the duodenum or ureters), or the presence of a contained rupture are critical observations. The type and size of the synthetic graft used, the location of the anastomoses, and the patency of the inferior mesenteric artery and lumbar arteries are also recorded in the operative report.

The pathology report on the resected aneurysm sac confirms the diagnosis of an abdominal aortic aneurysm, typically describing severe atherosclerosis with thinning and degeneration of the media, loss of elastic fibers, and chronic inflammation. It characterizes the wall tissue, often showing extensive lipid deposition, calcification, and ulceration. The report also rules out other pathologies such as inflammatory aneurysm (characterized by dense periaortic fibrosis and inflammatory cell infiltrates), mycotic (infected) aneurysm (presence of microorganisms), or connective tissue disorders (e.g., Marfan syndrome, Ehlers-Danlos syndrome, which show specific histological features of collagen or elastin defects), which might influence long-term management and surveillance. The presence of intraluminal thrombus is also confirmed.

\section{*{Expected Postoperative Course}}
A normal, successful postoperative course following open abdominal aortic aneurysm repair is characterized by a progressive and uneventful recovery, leading to a significant reduction in the risk of aneurysm rupture. Patients can expect initial discomfort and pain, which should be well-managed with analgesia and gradually decrease over days to weeks. Bowel function typically returns within 3-5 days, allowing for a gradual advancement of diet. The surgical incision should heal without signs of infection, excessive redness, swelling, or discharge. Mobility will gradually improve, with patients able to sit up, walk short distances, and perform basic self-care within the first few days, progressing to longer walks and light activities over several weeks.

Vital signs should remain stable, and urine output should be adequate, indicating good renal function. There should be no signs of major complications such as persistent bleeding, graft infection, bowel ischemia, or new neurological deficits. Patients will experience fatigue for several weeks to months, which is a normal part of recovery from major abdominal surgery. The ultimate goal is for the patient to return to their baseline functional status, free from aneurysm-related symptoms, with a durable repair that ensures long-term patency of the graft and restoration of normal blood flow to the lower extremities.

\section{*{Postoperative Care \& Follow-up}}
Long-term follow-up is essential after open abdominal aortic aneurysm repair to monitor for potential complications and manage cardiovascular risk factors.
\begin{itemize}
  \item **Postoperative Clinic Visits:** Initial follow-up visits are typically scheduled at 2-4 weeks, 3 months, and 6 months post-surgery to assess wound healing, discuss recovery progress, and address any concerns.
  \item **Suture/Staple Removal:** If non-absorbable skin sutures or staples were used, they are removed during an early follow-up visit, usually around 10-14 days post-op.
  \item **Imaging Surveillance:**
    \begin{itemize}
      \item An initial baseline computed tomography (CT) scan of the abdomen and pelvis with contrast is often performed 3-6 months post-surgery to confirm graft integrity, rule out early anastomotic pseudoaneurysms, and assess for any retroperitoneal fluid collections.
      \item Subsequent surveillance typically involves annual or biennial ultrasound or CT scans to monitor for anastomotic pseudoaneurysms, new aneurysms in other vascular beds (e.g., thoracic aorta, popliteal arteries), or graft-related issues such as kinking or stenosis.
    \end{itemize}
  \item **Cardiovascular Risk Factor Management:** Lifelong aggressive management of cardiovascular risk factors, including hypertension, hyperlipidemia, diabetes, and smoking cessation, is critical to prevent progression of atherosclerosis and reduce the risk of future cardiovascular events (e.g., myocardial infarction, stroke), which remain the leading cause of late mortality in these patients.
  \item **Physical Therapy/Rehabilitation:** Some patients, especially those with prolonged hospital stays or significant deconditioning, may benefit from a structured physical therapy or cardiac rehabilitation program to regain strength, mobility, and endurance.
  \item **Activity Resumption:** Patients are given guidelines for gradually resuming physical activities, typically avoiding heavy lifting (over 10-15 lbs) for several weeks to months to allow for complete abdominal wall healing and minimize the risk of incisional hernia.
  \item **Warning Signs:** Patients are educated on warning signs that require immediate medical attention, such as fever, severe or new abdominal pain, a pulsatile mass, signs of limb ischemia (e.g., new leg pain, numbness, pallor), or any signs of wound infection.
\end{itemize}
Adherence to these follow-up recommendations is crucial for optimizing long-term outcomes and detecting potential complications early.

\section*{Epidemiology}
Abdominal aortic aneurysms (AAAs) are a significant public health concern, with open repair being a critical intervention. The prevalence of AAA in the general population varies, but it is estimated to affect 4-8\% of men over 65 years of age and 0.5-1.5\% of women in the same age group. The incidence increases significantly with age, peaking in the seventh and eighth decades of life. Men are 4-6 times more likely to develop AAAs than women, although women tend to rupture at smaller aneurysm diameters and have a higher mortality rate from rupture. Major risk factors include smoking, advanced age, male sex, family history of AAA, hypertension, hyperlipidemia, and coronary artery disease. The primary indications for open repair include large aneurysm size (typically >5.5 cm in men, >5.0 cm in women), rapid expansion, or symptoms. While elective open AAA repair has a reported mortality rate of 2-5\% in experienced centers, the mortality rate for emergent open repair of a ruptured AAA remains exceedingly high, ranging from 50-80\%, underscoring the importance of timely elective intervention. Morbidity rates for elective open repair can be substantial, with complications such as cardiac events, renal failure, and respiratory complications occurring in 20-40\% of patients.

\section*{Outcomes \& Prognosis}
The outcomes and prognosis following elective open abdominal aortic aneurysm repair are generally excellent, particularly in carefully selected patients without severe comorbidities. For elective cases, the long-term survival rates are comparable to age-matched controls, provided that cardiovascular risk factors are aggressively managed. The 5-year survival rate for elective open AAA repair typically ranges from 70-85\%, with most late deaths attributable to cardiovascular events rather than aneurysm-related issues. Functional outcomes are also generally good, with most patients returning to their baseline activity levels within a few months. Recurrence of the aneurysm in the treated segment is not possible due to graft replacement, but patients remain at risk for developing new aneurysms in other locations (e.g., thoracic aorta, peripheral arteries) or anastomotic pseudoaneurysms, necessitating lifelong surveillance. Prognostic factors for favorable outcomes include younger age, fewer comorbidities (especially cardiac and renal), and elective rather than emergent presentation. Conversely, patients undergoing emergent repair for ruptured AAA face a significantly poorer prognosis, with high mortality rates despite successful surgery, primarily due to the physiological insult of hemorrhage and shock. Long-term success hinges on diligent follow-up and comprehensive management of systemic atherosclerosis, as highlighted by studies like the UK Small Aneurysm Trial and ADAM trial, which established the benefit of elective repair for larger aneurysms.

\section{*{References}}
\begin{enumerate}
  \item MedlinePlus. Abdominal aortic aneurysm repair - open. U.S. National Library of Medicine; 2024. Available from: \url{https://medlineplus.gov/ency/article/002956.htm}
  \item Sidawy AN, Perler BA, eds. Rutherford's Vascular Surgery and Endovascular Therapy. 10th ed. Elsevier; 2023.
  \item Sabiston DC, Townsend CM Jr, Beauchamp RD, Evers BM, Mattox KL, eds. Sabiston Textbook of Surgery: The Biological Basis of Modern Surgical Practice. 21st ed. Elsevier; 2021.
  \item Society for Vascular Surgery. Clinical Practice Guidelines for the Management of Abdominal Aortic Aneurysms. J Vasc Surg. 2018;67(1):2-77.e2.
\end{enumerate}

\clearpage